
\subsection{The Language and the Proofs}

§1 \quad So far, we have approached the idea of a class, and the relation of membership, largely through intuition. To develop set theory as a mathematical theory, however, we need to say more precisely what our language allows us to express and what counts as a legitimate inference. This also gives us an opportunity to explain a somewhat unusual feature of this textbook: alongside our ordinary mathematical proofs, we shall sometimes display their formal derivations in Gentzen's natural deduction calculus.

I do not intend to teach mathematical logic here. Readers unfamiliar with natural deduction should still be able to follow the mathematical development through the main text. Nevertheless, I want the formal foundations of our arguments to be available for inspection, rather than leaving them implicit. The accompanying Gentzen material therefore gives a complete specification of the language and inference rules we shall employ. Readers who wish to investigate these foundations can follow it in parallel; otherwise, the present explanation should suffice for understanding what the formal displays are doing.

§2 \quad We formulate NBG in a \textit{one-sorted first-order language with identity}. The expression \textit{one-sorted} means that all our variables range over objects of the same kind: classes. There are no separate variables for sets and proper classes, and no special quantifiers ranging exclusively over either. Instead, sethood will be introduced as a property that some classes have and others lack.

The language itself is remarkably economical. Its only primitive non-logical symbol is the membership relation $\in$, which we have already encountered. Identity, expressed by $=$, belongs to our logical vocabulary, together with the familiar connectives and quantifiers. Strictly speaking, our primitive language contains neither constants nor function symbols; its variables are its only terms. The various names, constructions, and other items of notation that we shall introduce as the theory develops must therefore be understood in terms of this underlying language.

\begin{formalderivation}
\subsubsection*{Formal Language of NBG}

\nbgpara{1}{Underlying Language}
We formulate NBG in a one-sorted first-order language with identity.

\nbgpara{2}{Variables}
The primitive variable symbols of the language are
\[
v_0,v_1,v_2,\ldots,
\]
indexed metalinguistically by the natural numbers. Thus, the language has denumerably infinitely many primitive variable symbols. Because the language contains neither constants nor function symbols, variables would be its only terms. We therefore dispense with a separate syntactic category of terms.
\nbgpara{3}{Display Names for Variables}
For readability, we shall normally display primitive variables using the customary letters
\[
x,y,z,a,b,c,\ldots
\qquad\text{and}\qquad
A,B,C,D,E,\ldots,
\]
adding subscripts whenever further variables are required. These customary letters are not additional variable symbols of the language. Rather, they are context-dependent typographical representatives of variables from the primitive stock
\[
v_0,v_1,v_2,\ldots.
\]

Within any formula, derivation, or continuous argument, each customary letter uniformly represents a particular primitive variable symbol, and distinct customary letters represent distinct primitive variable symbols. Thus, a displayed formula such as
\[
a\in A\to a\in B
\]
may, for example, represent the primitive formula
\[
v_{17}\in v_3\to v_{17}\in v_{26}.
\]
The choice of primitive variables represented by the customary letters may change between independent contexts. This display convention is entirely eliminable: every expression written using customary variable letters can be replaced uniformly by an expression using only variables from the primitive stock. 

\nbgpara{4}{Metavariables}
Throughout this text, we use
\[
\xi,\eta,\zeta,\ldots
\]
as metavariables ranging over primitive variable symbols of the language. Distinct metavariables are not assumed to represent distinct variables unless this is explicitly stated.

We use
\[
\varphi,\psi,\chi,\ldots
\]
as metavariables ranging over formulas of the language.

\nbgpara{5}{Logical Symbols}
The logical symbols are the identity symbol \(=\), the falsity symbol \(\bot\), the connectives
\[
\neg,\qquad
\land,\qquad
\lor,\qquad
\to,\qquad
\leftrightarrow,
\]
and the quantifiers
\[
\forall,\qquad
\exists.
\]
The symbols \(\lparen\) and \(\rparen\) are used as punctuation.

The logical vocabulary is deliberately not chosen for functional minimality. Although some of the connectives and one of the quantifiers could be defined in terms of the others, we take the familiar connectives and both quantifiers as primitive so that their respective inferential roles can be stated and studied directly. The non-logical vocabulary, by contrast, is chosen for maximal economy.

\nbgpara{6}{Non-Logical Symbol}
The sole non-logical primitive symbol is \(\in\), which is a binary relation symbol.

\nbgpara{7}{Atomic Formulas}
If \(\xi\) and \(\eta\) are variables, then
\[
\xi=\eta
\quad \text{ and } \quad
\xi\in\eta
\]
are atomic formulas.

An atomic formula of the form \(\xi\in\eta\) is read as ``\(\xi\) is a member of \(\eta\).''

\nbgpara{8}{Formation Rules}
The formulas of the language are generated recursively by the following clauses:

\begin{enumerate}
    \item Every atomic formula is a formula.

    \item \(\bot\) is a formula.

    \item If \(\varphi\) is a formula, then $\neg\varphi$ is a formula.

    \item If \(\varphi\) and \(\psi\) are formulas, then
    \[
    (\varphi\land\psi),\qquad
    (\varphi\lor\psi),\qquad
    (\varphi\to\psi),\qquad
    (\varphi\leftrightarrow\psi)
    \]
    are formulas.

    \item If \(\varphi\) is a formula and \(\xi\) is a variable, then
    \[
    \forall\xi\,\varphi
    \qquad\text{and}\qquad
    \exists\xi\,\varphi
    \]
    are formulas.
    \item Nothing else is a formula.
\end{enumerate}

\nbgpara{9}{Parenthesis Conventions}
Outermost parentheses may be omitted. Quantifiers and \(\neg\) bind more strongly than the binary connectives, while \(\land\) and \(\lor\) bind more strongly than \(\to\) and \(\leftrightarrow\). Connectives of equal binding strength have no precedence over one another.

Repeated conjunctions and disjunctions associate to the left. Thus,
\[
\varphi\land\psi\land\chi
\]
abbreviates
\[
(\varphi\land\psi)\land\chi,
\]
and analogously for \(\lor\). No associativity convention is adopted for \(\to\) or \(\leftrightarrow\).

The scope of a quantifier is the shortest formula immediately following it. Parentheses may be suppressed whenever these conventions determine a unique reading.

\nbgpara{10}{Free and Bound Variables}
An occurrence of a variable $\xi$ is bound if it lies within the scope of at least one quantifier $\forall \xi$ or $\exists \xi$. Otherwise, it is free. If an occurrence of \(\xi\) lies within the scopes of several quantifiers binding \(\xi\), it is governed by the innermost such quantifier.

A variable occurs free in a formula if at least one of its occurrences in that formula is free. A single variable may have both free and bound occurrences in the same formula.

\nbgpara{11}{Substitution}
Let \(\varphi\) be a formula, and let \(\xi\) and \(\eta\) be variables. We say that \(\eta\) is free for \(\xi\) in \(\varphi\) iff no free occurrence of \(\xi\) in \(\varphi\) lies within the scope of a quantifier binding \(\eta\).

If \(\eta\) is free for \(\xi\) in \(\varphi\), then
\[
\varphi[\eta/\xi]
\]
denotes the formula obtained by replacing every free occurrence of \(\xi\) in \(\varphi\) by \(\eta\). Bound occurrences of \(\xi\) are not replaced.

\nbgpara{12}{Alphabetic Variants}
Let \(Q\) be either \(\forall\) or \(\exists\). Suppose that \(\eta\) is distinct from \(\xi\) and does not occur in \(\varphi\). Then
\[
Q\xi\,\varphi
\]
and
\[
Q\eta\,\varphi[\eta/\xi]
\]

are called immediate alphabetic variants.

The substitution in the second formula replaces precisely those occurrences of \(\xi\) in \(\varphi\) that are bound by the displayed outer quantifier in the complete formula \(Q\xi\,\varphi\). Any occurrence of \(\xi\) already bound by a quantifier within \(\varphi\) is not replaced. A formula obtained from another by one or more successive renamings of this kind is called an alphabetic variant of that formula.

\nbgpara{13}{Intended Interpretation}
Our investigation of NBG will be syntactic: we shall study formulas, derivations, and the relations of derivability that hold between them. Nevertheless, the intended interpretation may serve as an informal guide to the formalism.

Under the intended interpretation of NBG, the variables range over classes. Since the language is one-sorted, every variable ranges over the same domain. The distinction between sets and proper classes is therefore not built into the variable symbols or the quantifiers, but will instead be expressed within the language.

Distinct variable symbols need not be understood as representing distinct classes. Thus, although \(v_0\) and \(v_{15}\) are distinct variable symbols, their metalinguistic distinctness does not by itself license the object-language formula
\[
\neg(v_0= v_{15}).
\]
Intuitively, identity and non-identity are relations between the objects represented by variables, not relations between the variable symbols themselves.

\nbgpara{14}{Typographical Convention}
We shall generally use lowercase variables when emphasizing an object's role as a member and capital variables when emphasizing its role as a class whose members are under consideration. Thus, in
\[
a\in A,
\]
the notation emphasizes the role of \(a\) as a member and the role of \(A\) as a class whose members are under consideration.

This convention has no formal significance. Lowercase and capital variables are variables of exactly the same kind, range over the same domain, and may occur in the same syntactic positions. In particular, lowercase variables do not range only over sets, and capital variables do not range only over proper classes. 



\end{formalderivation}

§3 \quad A word about our variables. We shall generally use lowercase letters when emphasizing an object's role as a member, and capital letters when emphasizing its role as a class whose members are under consideration. Thus, in $a\in A$, we think of $a$ as a member of $A$. This is only a typographical convention. Lowercase letters do not range only over sets, nor do capital letters range only over proper classes. Both range over classes without restriction, and either may occur in either position of the membership relation.

There is a further distinction that is easy to overlook: different variable symbols need not represent different objects. That $a$ and $b$ are different letters does not establish that $a\ne b$. Identity concerns the objects under consideration, not the symbols with which we represent them.

§4 \quad Our proofs will ordinarily be explained in words, much as in other mathematics textbooks. But selected proofs will also have a more explicit counterpart in the Gentzen sidebar. In a Gentzen derivation, the formulas are arranged in a tree, which we read from top to bottom. Horizontal lines represent inferences: the formulas above a line license the conclusion below it, according to the rule indicated on the right.

For a simple example, consider the following derivation:

\begin{formalderivation}
\begin{prooftree}
    \AxiomC{}
    \RightLabel{\(1\)}
    \UnaryInfC{\(a\in A\land a\in B\)}
    \RightLabel{\(\land\)-E}
    \UnaryInfC{\(a\in A\)}
    \RightLabel{\(\to\)-I, \(1\)}
    \UnaryInfC{\((a\in A\land a\in B)\to a\in A\)}
\end{prooftree}
\end{formalderivation}

We begin by assuming that $a$ belongs to both $A$ and $B$. The first inference, labelled $\land$-E, expresses the familiar fact that from a conjunction we may infer either of its conjuncts. We therefore obtain $a\in A$. Having established this on the assumption that $a\in A\land a\in B$, we may conclude that the latter implies the former. This is an application of implication introduction ($\to$-I).

§5 \quad Notice the number $1$ beside the initial assumption and again beside the final inference. The number tells us that the assumption is \textit{discharged} by that application of $\to$-I. In other words, we used the assumption in the course of our reasoning, but the conclusion is no longer dependent on its being true. The derivation has no undischarged assumptions and is therefore a \textit{proof}. Its conclusion is a theorem of our logical system.

This distinction is worth keeping in mind. A derivation may establish a conclusion only on certain assumptions, whereas a proof establishes its conclusion without any undischarged assumptions. The Gentzen material makes this precise by defining the dependency of each formula occurrence on assumptions. Assumption labels are merely bookkeeping devices for recording when those dependencies are discharged; they are not part of the mathematical formulas themselves.

\begin{formalderivation}
\subsubsection*{Logical Inference Rules: Derivations and Dependency}

\nbgpara{1}{Derivations}
A derivation is a finite rooted tree of formula occurrences constructed in accordance with the inference rules stated below. Its lowermost formula occurrence is its conclusion. Each uppermost formula occurrence is either an assumption or the conclusion of a rule with no premises.

Since derivations are trees of formula occurrences, distinct nodes may contain the same formula while occupying different positions in the derivation and depending on different assumptions.

\nbgpara{2}{Assumptions and Dependency}
An assumption occurrence may be marked with an assumption label. The assumption labels are natural numbers. We use \(n,m,\ldots\) as metavariables ranging over assumption labels. Thus, if \(n\) is an assumption label, then \(n+1\) is also an assumption label. An assumption occurrence marked with the label \(n\) is written as follows:

\begin{prooftree}
    \AxiomC{}
    \RightLabel{\(n\)}
    \UnaryInfC{\(\varphi\)}
\end{prooftree}

Assumption labels are metalinguistic bookkeeping marks. They are not symbols of the language of NBG and are not parts of the formulas to which they are attached.

Within a derivation, each assumption label identifies a single discharge class. All assumption occurrences carrying the same label are occurrences of the same formula, and a label is associated with at most one application of a discharge rule. If a label is associated with such an application, every occurrence carrying that label lies within the premise subderivation designated by the rule. The same labels may be reused in separate derivations.

When an application of a discharge rule carries the label \(n\), it discharges every still-undischarged assumption occurrence marked with \(n\) in the premise subderivation designated by that rule. A single application may therefore discharge several occurrences of the same assumption.

For \(\to\)-I, \(\neg\)-I, and CR, the designated premise subderivation is the sole premise subderivation of the rule. In an application of \(\lor\)-E, each label applies only to its corresponding minor premise subderivation. In an application of \(\exists\)-E, the label applies only to the minor premise subderivation beginning with the assumed instance of the existential formula.

The dependency of formula occurrences on assumptions is defined recursively. For each formula occurrence \(\theta\) in a derivation, let \(\Dep(\theta)\) be its dependency set.

\begin{enumerate}
    \item If \(\theta\) is an assumption occurrence, then
    \[
    \Dep(\theta)=\{\theta\}.
    \]

    \item If \(\theta\) is the conclusion of a rule with no premises, then
    \[
    \Dep(\theta)=\varnothing.
    \]

    \item Suppose that \(\theta\) is the conclusion of an application of a rule whose premise occurrences are
    \[
    \theta_1,\ldots,\theta_k.
    \]
    Let \(\Delta\) be the set of assumption occurrences discharged by that application. Then
    \[
    \Dep(\theta)
    =
    \bigl(
    \Dep(\theta_1)\cup\cdots\cup\Dep(\theta_k)
    \bigr)
    \setminus\Delta.
    \]
\end{enumerate}

A formula occurrence \(\theta\) depends on precisely the assumption occurrences belonging to \(\Dep(\theta)\). The undischarged assumptions of a derivation are the assumption occurrences belonging to the dependency set of its conclusion.

\nbgpara{3}{Turnstile}
Let \(\Gamma\) be a collection of formulas and let \(\varphi\) be a formula. We write
\[
\Gamma\vdash\varphi
\]
iff there exists a derivation of \(\varphi\) such that the formula occurring at every undischarged assumption occurrence is a member of \(\Gamma\). This is read as ``\(\varphi\) is derivable from \(\Gamma\).'' Thus, the symbol \(\vdash\) is a symbol of the metalanguage, not of the language of NBG.

The definition does not require every member of \(\Gamma\) to occur as an undischarged assumption. Consequently, if
\[
\Gamma\vdash\varphi
\qquad\text{and}\qquad
\Gamma\subseteq\Delta,
\]
then
\[
\Delta\vdash\varphi.
\]
Weakening is therefore built into the definition of derivability.

The collection \(\Gamma\) also does not record the multiplicity of assumption occurrences. Several undischarged occurrences of the same formula are represented by a single member of \(\Gamma\), although they remain distinct occurrences within the derivation itself.

\nbgpara{4}{Proofs}
A derivation with no undischarged assumptions is called a proof. Its conclusion is called a theorem of the deductive system in which the derivation is carried out.

When the derivation has no undischarged assumptions, we write
\[
\vdash\varphi
\]
instead of
\[
\varnothing\vdash\varphi.
\]
Thus, \(\vdash\varphi\) means that there exists a proof of \(\varphi\), or equivalently that \(\varphi\) is a theorem of the deductive system.

\vspace{0.5cm}


\end{formalderivation}

§6 \quad Each logical operation is governed by corresponding inference rules. The rules of conjunction and implication, for example, license the two steps in our illustration. Other rules govern disjunction, negation, the quantifiers, and identity. We also include classical reductio, allowing us to establish a statement by showing that its negation leads to a contradiction.

The rules have names, usually formed from the logical symbol together with \textit{I} for introduction or \textit{E} for elimination. In a displayed proof, these names appear beside the corresponding inference lines. Certain applications, especially those involving quantifiers, must also satisfy additional conditions. The complete rules and their conditions are given in the sidebar; the ordinary proofs will explain the relevant mathematical reasoning without requiring the reader to memorize them in advance.

\begin{formalderivation}
\subsubsection*{Propositional Inference Rules}

\nbgpara{5}{Conjunction}

\nbgparar{$\bm{\land}$-I} \quad Conjunction is governed by the following introduction rule:

\begin{prooftree}
    \AxiomC{\(\varphi\)}
    \AxiomC{\(\psi\)}
    \RightLabel{\(\land\)-I}
    \BinaryInfC{\(\varphi\land\psi\)}
\end{prooftree}

\vspace{0.3cm}

\nbgparar{$\bm{\land}$-E} \quad Conjunction is governed by the following two elimination rules:

\begin{center}
\begin{minipage}{0.55\textwidth}
\centering
\begin{prooftree}
    \AxiomC{\(\varphi\land\psi\)}
    \RightLabel{\(\land\)-E}
    \UnaryInfC{\(\varphi\)}
\end{prooftree}
\end{minipage}
\hfill
\begin{minipage}{0.42\textwidth}
\centering
\begin{prooftree}
    \AxiomC{\(\varphi\land\psi\)}
    \RightLabel{\(\land\)-E}
    \UnaryInfC{\(\psi\)}
\end{prooftree}
\end{minipage}
\end{center}

\vspace{0.5cm}

\nbgpara{6}{Disjunction}

\nbgparar{$\bm{\lor}$-I} \quad Disjunction is governed by the following two introduction rules:

\begin{center}
\begin{minipage}{0.55\textwidth}
\centering
\begin{prooftree}
    \AxiomC{\(\varphi\)}
    \RightLabel{\(\lor\)-I}
    \UnaryInfC{\(\varphi\lor\psi\)}
\end{prooftree}
\end{minipage}
\hfill
\begin{minipage}{0.42\textwidth}
\centering
\begin{prooftree}
    \AxiomC{\(\psi\)}
    \RightLabel{\(\lor\)-I}
    \UnaryInfC{\(\varphi\lor\psi\)}
\end{prooftree}
\end{minipage}
\end{center}

\vspace{0.3cm}

\nbgparar{$\bm{\lor}$-E} \quad Disjunction is governed by the following elimination rule:

\begin{prooftree}
    \AxiomC{\(\varphi\lor\psi\)}

    \AxiomC{}
    \RightLabel{\(n\)}
    \UnaryInfC{\(\varphi\)}
    \noLine
    \UnaryInfC{\(\vdots\)}
    \noLine
    \UnaryInfC{\(\chi\)}

    \AxiomC{}
    \RightLabel{\(m\)}
    \UnaryInfC{\(\psi\)}
    \noLine
    \UnaryInfC{\(\vdots\)}
    \noLine
    \UnaryInfC{\(\chi\)}

    \RightLabel{\(\lor\)-E, \(n,m\)}
    \TrinaryInfC{\(\chi\)}
\end{prooftree} \vspace{0.3cm}

The labels \(n\) and \(m\) are distinct. The application discharges the assumption occurrences marked with \(n\) only in the left minor premise subderivation and the assumption occurrences marked with \(m\) only in the right minor premise subderivation.

\vspace{0.3cm}

\nbgpara{7}{Implication}

\nbgparar{$\bm{\to}$-I} \quad Implication is governed by the following introduction rule:

\begin{prooftree}
    \AxiomC{}
    \RightLabel{\(n\)}
    \UnaryInfC{\(\varphi\)}
    \noLine
    \UnaryInfC{\(\vdots\)}
    \noLine
    \UnaryInfC{\(\psi\)}
    \RightLabel{\(\to\)-I, \(n\)}
    \UnaryInfC{\(\varphi\to\psi\)}
\end{prooftree}

The application discharges every assumption occurrence marked with \(n\) in its premise subderivation.

An application of \(\to\)-I may be vacuous: the rule may be applied even if its premise subderivation contains no assumption occurrence marked with the indicated label \(n\):

\begin{prooftree}
    \AxiomC{}
    \noLine
    \UnaryInfC{\(\vdots\)}
    \noLine
    \UnaryInfC{\(\psi\)}
    \RightLabel{\(\to\)-I, \(n\)}
    \UnaryInfC{\(\varphi\to\psi\)}
\end{prooftree}

\vspace{0.3cm}

\nbgparar{$\bm{\to}$-E} \quad Implication is governed by the following elimination rule:

\begin{prooftree}
    \AxiomC{\(\varphi\to\psi\)}
    \AxiomC{\(\varphi\)}
    \RightLabel{\(\to\)-E}
    \BinaryInfC{\(\psi\)}
\end{prooftree}

\nbgpara{8}{Biconditional}

\nbgparar{$\bm{\leftrightarrow}$-I} \quad The biconditional is governed by the following introduction rule:

\begin{prooftree}
    \AxiomC{\(\varphi\to\psi\)}
    \AxiomC{\(\psi\to\varphi\)}
    \RightLabel{\(\leftrightarrow\)-I}
    \BinaryInfC{\(\varphi\leftrightarrow\psi\)}
\end{prooftree}

\vspace{0.3cm}

\nbgparar{$\bm{\leftrightarrow}$-E} \quad The biconditional is governed by the following two elimination rules:

\begin{center}
\begin{minipage}{0.55\textwidth}
\centering
\begin{prooftree}
    \AxiomC{\(\varphi\leftrightarrow\psi\)}
    \AxiomC{\(\varphi\)}
    \RightLabel{\(\leftrightarrow\)-E}
    \BinaryInfC{\(\psi\)}
\end{prooftree}
\end{minipage}
\hfill
\begin{minipage}{0.42\textwidth}
\centering
\begin{prooftree}
    \AxiomC{\(\varphi\leftrightarrow\psi\)}
    \AxiomC{\(\psi\)}
    \RightLabel{\(\leftrightarrow\)-E}
    \BinaryInfC{\(\varphi\)}
\end{prooftree}
\end{minipage}
\end{center} 

\nbgpara{9}{Negation}

\nbgparar{$\bm{\neg}$-I} \quad Negation is governed by the following introduction rule:

\begin{prooftree}
    \AxiomC{}
    \RightLabel{\(n\)}
    \UnaryInfC{\(\varphi\)}
    \noLine
    \UnaryInfC{\(\vdots\)}
    \noLine
    \UnaryInfC{\(\bot\)}
    \RightLabel{\(\neg\)-I, \(n\)}
    \UnaryInfC{\(\neg\varphi\)}
\end{prooftree}

\vspace{0.3cm}

\nbgparar{$\bm{\neg}$-E} \quad Negation is governed by the following elimination rule:

\begin{prooftree}
    \AxiomC{\(\varphi\)}
    \AxiomC{\(\neg\varphi\)}
    \RightLabel{\(\neg\)-E}
    \BinaryInfC{\(\bot\)}
\end{prooftree}


\nbgpara{10}{Falsity}

\nbgparar{$\bm{\bot}$-E} \quad From falsity, any formula may be inferred:

\begin{prooftree}
    \AxiomC{\(\bot\)}
    \RightLabel{\(\bot\)-E}
    \UnaryInfC{\(\varphi\)}
\end{prooftree}

\vspace{0.3cm}

\nbgpara{11}{Classical Reductio}

\nbgparar{CR} \quad Classical reductio is governed by the following rule:

\begin{prooftree}
    \AxiomC{}
    \RightLabel{\(n\)}
    \UnaryInfC{\(\neg\varphi\)}
    \noLine
    \UnaryInfC{\(\vdots\)}
    \noLine
    \UnaryInfC{\(\bot\)}
    \RightLabel{CR, \(n\)}
    \UnaryInfC{\(\varphi\)}
\end{prooftree} \vspace{0.3cm}

The propositional rules stated in paragraphs 5--10 are intuitionistically valid. The addition of CR is what makes the propositional fragment of the deductive system classical. \vspace{-0.3cm}


\end{formalderivation}

\begin{formalderivation}
\subsubsection*{Quantification, Identity, and Display Conventions}

\nbgpara{12}{Universal Quantification}

\nbgparar{$\bm{\forall}$-I} \quad Universal quantification is governed by the following introduction rule: \vspace{0.3cm}

\begin{prooftree}
    \AxiomC{\(\varphi[\eta/\xi]\)}
    \RightLabel{\(\forall\)-I}
    \UnaryInfC{\(\forall\xi\,\varphi\)}
\end{prooftree} \vspace{0.3cm}

The variable $\eta$ is called the eigenvariable of the application. Intuitively, it is treated as arbitrary: no undischarged assumption on which the premise depends may contain $\eta$ free.

This is made precise by the following side conditions:

\begin{enumerate}
    \item \(\eta\) is free for \(\xi\) in \(\varphi\);

    \item \(\eta\) does not occur free in \(\varphi\); and

    \item \(\eta\) does not occur free in any undischarged assumption on which the premise occurrence \(\varphi[\eta/\xi]\) depends.
\end{enumerate}

\nbgparar{$\bm{\forall}$-E} \quad Universal quantification is governed by the following elimination rule: \vspace{0.3cm}

\begin{prooftree}
    \AxiomC{\(\forall\xi\,\varphi\)}
    \RightLabel{\(\forall\)-E}
    \UnaryInfC{\(\varphi[\eta/\xi]\)}
\end{prooftree} \vspace{0.3cm}

The rule is permitted only if \(\eta\) is free for \(\xi\) in \(\varphi\).

\vspace{0.5cm}

\nbgpara{13}{Existential Quantification}

\nbgparar{$\bm{\exists}$-I} \quad Existential quantification is governed by the following introduction rule: \vspace{0.3cm}

\begin{prooftree}
    \AxiomC{\(\varphi[\eta/\xi]\)}
    \RightLabel{\(\exists\)-I}
    \UnaryInfC{\(\exists\xi\,\varphi\)}
\end{prooftree} \vspace{0.3cm}

The rule is permitted only if \(\eta\) is free for \(\xi\) in \(\varphi\).

\vspace{0.3cm}

\nbgparar{$\bm{\exists}$-E} \quad Existential quantification is governed by the following elimination rule:

\begin{prooftree}
    \AxiomC{\(\exists\xi\,\varphi\)}

    \AxiomC{}
    \RightLabel{\(n\)}
    \UnaryInfC{\(\varphi[\eta/\xi]\)}
    \noLine
    \UnaryInfC{\(\vdots\)}
    \noLine
    \UnaryInfC{\(\psi\)}

    \RightLabel{\(\exists\)-E, \(n\)}
    \BinaryInfC{\(\psi\)}
\end{prooftree} \vspace{0.3cm}

The right-hand premise subderivation ending with the displayed occurrence of $\psi$ is called the minor premise subderivation. It contains the assumption occurrences $\varphi[\eta/\xi]$ marked with $n$, which are discharged by the application. The application of \(\exists\)-E discharges the assumption occurrences marked with \(n\) only within this subderivation.

The variable \(\eta\) is called the eigenvariable of the application. Its role is to represent an arbitrary witness to the existential statement about which no other relevant assumption has been made.

This is made precise by the following side conditions:

\begin{enumerate}
    \item \(\eta\) is free for \(\xi\) in \(\varphi\);

    \item \(\eta\) does not occur free in \(\exists\xi\,\varphi\);

    \item \(\eta\) does not occur free in \(\psi\); and

    \item \(\eta\) does not occur free in any undischarged assumption on which the minor premise derivation of \(\psi\) depends, other than the assumption occurrences discharged by the application.
\end{enumerate}

\vspace{0.35cm}

\nbgpara{14}{Identity}

\nbgparar{$\bm{=}$-I} \quad Identity is governed by the following introduction rule:

\begin{prooftree}
    \AxiomC{}
    \RightLabel{\(=\)-I}
    \UnaryInfC{\(\xi=\xi\)}
\end{prooftree}

\vspace{0.2cm}

\nbgparar{$\bm{=}$-E} \quad Identity is governed by the following two elimination rules: \vspace{0.3cm}

\begin{center}
\begin{minipage}{0.5\textwidth}
\centering
\begin{prooftree}
    \AxiomC{\(\xi=\eta\)}
    \AxiomC{\(\varphi[\xi/\zeta]\)}
    \RightLabel{\(=\)-E}
    \BinaryInfC{\(\varphi[\eta/\zeta]\)}
\end{prooftree}
\end{minipage}
\hfill
\begin{minipage}{0.46\textwidth}
\centering
\begin{prooftree}
    \AxiomC{\(\xi=\eta\)}
    \AxiomC{\(\varphi[\eta/\zeta]\)}
    \RightLabel{\(=\)-E}
    \BinaryInfC{\(\varphi[\xi/\zeta]\)}
\end{prooftree}
\end{minipage}
\end{center} \vspace{0.3cm}

Each rule is permitted only if both \(\xi\) and \(\eta\) are free for \(\zeta\) in \(\varphi\).


\nbgpara{15}{Compressed Derivation Displays}
In later derivations, we shall occasionally suppress a finite sequence of routine applications of the same inference rule. If \(R\) is an inference rule, a displayed application labelled

$$
R^{*}
$$

indicates that the displayed premise occurrences and conclusion are connected by a finite nonempty derivation segment whose inference applications are all applications of \(R\). Intermediate formula occurrences within that segment are omitted from the display.

Thus, for example, a step of the form

\begin{prooftree}
    \AxiomC{\(\varphi\)}
    \RightLabel{\(\forall\)-I\(^{*}\)}
    \UnaryInfC{\(\forall\xi_1\cdots\forall\xi_k\,\varphi\)}
\end{prooftree}

abbreviates the corresponding succession of applications of \(\forall\)-I, with the quantifiers introduced one at a time in the order required to obtain the displayed conclusion.

Similarly, a displayed application of \(\land\)-I\(^{*}\), \(\land\)-E\(^{*}\), or another rule may abbreviate the corresponding finite derivation segment.

This convention concerns only the display of derivations and introduces no new inference rules. Every suppressed application must be a legitimate application of the ordinary rule and must separately satisfy all side conditions governing that rule.

\end{formalderivation}

§7 \quad The formal derivations are intended to supplement the main exposition, not to replace it. A tree can tell us exactly which inference rules have been applied, but it does not necessarily tell us why someone would think of making those inferences, or what the resulting theorem means. These are questions for the ordinary mathematical discussion. Conversely, a prose proof may move rather quickly over steps whose logical structure is worth inspecting in detail. The two presentations therefore have different, but complementary, purposes.

As the theory develops, we shall also introduce \textit{derived rules}: patterns of inference that have been justified by means of the rules already available. Such rules are more than convenient abbreviations. Once we have understood a pattern of reasoning, we can often make the same inference intuitively, without working through its individual steps anew. This has a precise formal counterpart: each application of a derived rule can be replaced by an appropriate derivation using rules already established. We shall return to this idea with our first example, the rule Set-I.

§8 \quad There is a final point about the status of this formal development. The Gentzen trees are not merely illustrations intended to make our arguments look rigorous. They belong to a precisely specified deductive system. In addition to stating its primitive inference rules, I have developed the metatheoretical results needed to justify such operations as combining derivations, citing previously established theorems, and introducing derived rules. The technical material is therefore not an alternative set of assumptions about classes; it explains how the reasoning by which we develop those assumptions and their consequences is formally organized.

As we enrich the language with definitions and, eventually, notation for constructing classes, further justification will be required. In particular, it matters that convenient notation can be eliminated in favor of expressions of the primitive language, and that the rules governing its use are legitimate. The corresponding technical development will accompany the mathematics where it becomes relevant. For now, we have established the language in which NBG will be expressed and the logical rules by which its arguments will proceed.
